% Generated draft: requires mathematical and source review.
\documentclass[11pt]{article}
\usepackage[margin=25mm]{geometry}
\usepackage{fontspec}
\setmainfont{Latin Modern Roman}
\usepackage{amsmath,amssymb,mathtools,unicode-math}
\setmathfont{Latin Modern Math}
\usepackage{enumitem,needspace,xurl,hyperref}
\setlength{\parindent}{0pt}
\setlength{\parskip}{5pt}
\setlength{\emergencystretch}{3em}
\newcommand{\N}{\mathbb{N}}
\newcommand{\Z}{\mathbb{Z}}
\newcommand{\Q}{\mathbb{Q}}
\newcommand{\R}{\mathbb{R}}
\newcommand{\C}{\mathbb{C}}
\newcommand{\F}{\mathbb{F}}
\newcommand{\A}{\mathbb{A}}
\newcommand{\PP}{\mathbb{P}}
\newcommand{\E}{\mathbb{E}}
\newcommand{\eps}{\varepsilon}
\newcommand{\dd}{\,\mathrm{d}}
\begin{document}
\section*{UnsolvedMath ID: 30004569}
\textbf{Source problem code:} catalogue entry
\subsection{Definitions and mathematical statement}
\textbf{Scope and terminology.} The supplied record concerns a many-particle kinetic model and the quasilinear approximation of its Vlasov equation. A deterministic regime refers to evolution determined by the prescribed initial data and dynamics; a stochastic regime refers to an effective description involving randomness. The source describes the long-time stochastic regime as one in which randomness drives effective diffusion. It does not specify the kinetic equation, the quasilinear approximation, the source or probability law of the randomness, or the effective diffusion operator.

\textbf{Statement of the source question.} Can one rigorously characterize the transition, in the quasilinear approximation of the Vlasov equation for the many-particle model discussed in the cited report, from a short-time deterministic regime to a long-time stochastic regime in which randomness drives effective diffusion?

\textbf{Unresolved mathematical qualifications.} This is a research question rather than a fully specified conjecture. The supplied record gives no phase-space domain, interaction law, initial-data class, boundary conditions, approximation parameter, time scaling, transition criterion, or mode of convergence. Consequently, the terms ``short-time,'' ``long-time,'' and ``effective diffusion'' cannot be assigned quantitative meanings from this record, and no precise quantified limit statement can be supplied without additional hypotheses. The original formulation refers generally to the many-particle kinetic model in the report; the supplied corrected formulation specifically identifies the quasilinear Vlasov approximation.
\subsection{Short English statement}
The problem asks for a rigorous description of how the quasilinear Vlasov approximation passes from deterministic behavior at short times to stochastic behavior with randomness-driven effective diffusion at long times. The supplied record leaves the model, scaling, and meaning of this transition unspecified, so a precise theorem statement requires further source information.
\subsection{Sources}
\begin{itemize}
\item Local registry: \texttt{data/unsolvedmath/problems.json}, permanent UnsolvedMath ID \(30004569\), exact source problem code \texttt{catalogue entry}; supplied title: \emph{Deterministic-to-Stochastic Transition in the Vlasov Approximation}.
\item Supplied report citation: \emph{Classical and Quantum Mechanical Models of Many-Particle Systems} (2021), Oberwolfach Reports, pp.\ 1857--1902, DOI \texttt{10.4171/owr/2020/38}. Supplied source and literature link: \url{https://doi.org/10.4171/owr/2020/38}. The record describes Oberwolfach Reports as documenting mathematical workshops and their open-problem discussions.
\item Additional supplied literature reference: ``Bardos, deterministic and stochastic Vlasov dynamics,'' \url{https://www.cirm-math.fr/RepOrga/2355/Abstracts/Claude_BARDOS.pdf}. This reference does not establish attribution of the problem to Bardos; the record names no proposer.
\item Supplied literature assessment, checked on exactly \texttt{2026-08-22}: the record labels the problem open and describes the deterministic-to-stochastic transition for Vlasov-type dynamics as a recognized challenging problem, with available mean-field and fluctuation results not providing the requested general transition theory. This is supplied dated curation derived from the OWR curation, not an independently verified literature assessment; its time-sensitive status and cited claims require independent verification.
\item The supplied record labels its statement \texttt{corrected\_verified} and reports that the official Oberwolfach report was checked, the question isolated from surrounding exposition, and notation repaired for self-contained publication. These are supplied verification claims, not verification performed for this draft. The proposed year and report-year tag are 2020, whereas the supplied bibliographic citation gives 2021.
\end{itemize}
\end{document}
